\subsection{Invariant means}

Let X be a topological space. Let \(\mathcal{B}(\mathrm{X};\mathbf{R})\) denote the real vector space consisting of continuous bounded mappings from X into R. Endowed with the norm \(\| f\| = \sup_{x\in \mathbf{X}}|f(x)|\) , this is a Banach space (GT, X, § 3, No.~1); it is also an ordered vector space, where the relation \(f\geqslant g\) means \(\ll f(x)\geqslant g(x) \text{ for all } x\in \mathbf{X}\gg\) .

DEFINITION 1. --- A positive linear form \(\mu\) on the space \(\mathcal{B}(\mathbf{X};\mathbf{R})\) , where \(\mathbf{X}\) is a topological space, for which \(\|\mu\| = 1\) , is called a mean on \(\mathbf{X}\) .

* When X is compact, a mean on X is a positive measure on X such that \(\mu(X)=1\) . *

Lemma 1. --- The set K of means on X is the subset of the unit ball of the dual of the Banach space \(E = \mathcal{B}(X; \mathbf{R})\) whose elements are the linear forms \(\mu\) such that \(\mu(1)=1\) . It is a subset of \(E'\) which is convex and compact for \(\sigma(E', E)\) .

Let \(\mu\) be a linear form on \(E\) , such that \(\mu(1) = 1\) . For every function \(f \in E\) , we define the function \(f' \in E\) by \(f'(x) = \| f \| - f(x) (x \in X)\) . First assume that \(\mu\) is a mean; for every \(f \in E\) , we have \(f' \geqslant 0\) , hence \(\mu(f') \geqslant 0\) , i.e.~\(\mu(f) \leqslant \| f \|\) ; therefore \(\|\mu\| \leqslant 1\) .

Conversely, suppose \(\mu\) belongs to \(E'\) , and that \(\| \mu \| \leqslant 1\) ; for every positive function \(f \in E\) , we have \(\mu(f') \leqslant \| f'\|\) , hence

\[
\| f \| - \mu (f) = \mu \left(f ^ {\prime}\right) \leqslant \| f ^ {\prime} \| \leqslant \| f \|,
\]

and finally \(\mu(f) \geqslant 0\) ; consequently, \(\mu\) is a mean.

It is clear that K is convex; that it is compact for \(\sigma(\mathrm{E}^{\prime},\mathrm{E})\) follows from cor. 3 of III, p.~17.

Q.E.D.

Let \(\Gamma\) be a set of continuous mappings from X into X which commute pairwise. Let \(\gamma \in \Gamma\) . For every function \(f \in \mathbf{E}\) , we have \(f \circ \gamma \in \mathbf{E}\) ; hence we can define an affine transformation \(u_{\gamma}\) on the set K of means on X, by

\[
u _ {\gamma} \mu (f) = \mu (f \circ \gamma) (\mu \in \mathrm{K}, f \in \mathrm{E}).
\]

If K is assigned the topology induced by \(\sigma(\mathrm{E}^{\prime},\mathrm{E})\) , the mapping \(u_{\gamma}\) is continuous. If \(\gamma\) is a homeomorphism, \(u_{\gamma}\mu\) can be deduced from \(\mu\) by transport of structure. Finally, we have \(u_{\gamma}u_{\gamma^{\prime}} = u_{\gamma^{\prime}}u_{\gamma}\) for all \(\gamma\) , \(\gamma^{\prime}\) in \(\Gamma\) . By the Markoff-Kakutani th. (IV, p.~39, th. 1), there exists a mean \(\mu\) on X, such that \(u_{\gamma}\mu = \mu\) for all \(\gamma \in \Gamma\) ; in other words, \(\mu\) satisfies the relation \(\mu(f) = \mu(f \circ \gamma)\) for \(f \in E\) and \(\gamma \in \Gamma\) .

In an analogous way, the corollary of th. 1 (IV, p.~40) implies the following result :

PROPOSITION 1. --- Let X be a topological space and G a solvable group. We assume that G operates on X on the left, in such a way that for all \(g \in G\) , the mapping \(x \mapsto g \cdot x\) from X into X is continuous. Then there exists a mean on X which is invariant under G.

COROLLARY. --- Let G be a solvable topological group. Then there exists a mean on G which is invariant under the left and the right translations.

It is enough to apply prop. 1 to the solvable group \(\mathbf{G} \times \mathbf{G}\) acting on \(\mathbf{G}\) by \((g, g') \cdot x = gxg'^{-1}\) .

